\section{Custody: the Taproot tree that plays the game}\label{sec:custody}

\subsection{The state output}

A frog in custody is a Taproot output with the fixed internal key
\code{NUMS_XONLY}:
\begin{center}
\footnotesize\code{50929b74c1a04954b78b4b6035e97a5e078a5a0f28ec96d547bfee9ace803ac0}.
\end{center}
No discrete log of this point is known, so there is no key-path spend. The
output can only be spent through one of four leaves. The tree is balanced at
depth two: three crossing leaves $L,S,R$ (left, straight, right) and one send
leaf $X$ (transfer).

\begin{align}
h_L&=\mathrm{leaf\_hash}(L), & h_S&=\mathrm{leaf\_hash}(S),\notag\\
h_R&=\mathrm{leaf\_hash}(R), & h_X&=\mathrm{leaf\_hash}(X),\notag\\
h_{LS}&=\mathrm{branch}(h_L,h_S), & h_{RX}&=\mathrm{branch}(h_R,h_X),\notag\\
\mathit{root}&=\mathrm{branch}(h_{LS},h_{RX}), & Q&=\mathrm{taproot\_output\_key}(\text{NUMS},\mathit{root}),\notag
\end{align}
where
\begin{equation}
\mathrm{branch}(a,b)=\mathrm{tagged\_hash}(\text{``TapBranch''},
  \min(a,b)\,\|\,\max(a,b)),
\end{equation}
and
\begin{equation}
\mathrm{leaf\_hash}(s)=\mathrm{tagged\_hash}(\text{``TapLeaf''},
  \text{\code{0xc0}}\,\|\,\mathrm{compact\_size}(|s|)\,\|\,s).
\end{equation}

\begin{figure}[htbp]
\centering
\begin{tikzpicture}[
  font=\small,
  leaf/.style={draw=deepgreen, fill=pale, rectangle, align=center,
               text width=23mm, minimum height=11mm, inner sep=3pt,
               font=\scriptsize\sffamily},
  branch/.style={draw=midgreen, fill=white, rectangle, align=center,
                 inner sep=4pt, font=\scriptsize\sffamily},
  root/.style={draw=deepgreen, fill=deepgreen!8, rectangle, align=center,
               inner sep=5pt, font=\scriptsize\sffamily, text width=56mm},
  edge/.style={draw=deepgreen!70, line width=0.7pt},
  lbl/.style={font=\tiny\sffamily, color=muted, fill=white, inner sep=1pt}
]
\node[root] (root) at (0,3)
  {output key $Q$\\\code{taproot_output_key(NUMS, root)}};
\node[branch] (ls) at (-4.4,1) {$h_{LS}=\mathrm{branch}(h_L,h_S)$};
\node[branch] (rx) at (4.4,1) {$h_{RX}=\mathrm{branch}(h_R,h_X)$};
\node[leaf] (L) at (-6.7,-1.3) {\textbf{$L$ — left}\\body dir $=0$};
\node[leaf] (S) at (-2.1,-1.3) {\textbf{$S$ — straight}\\body dir $=1$};
\node[leaf] (R) at (2.1,-1.3) {\textbf{$R$ — right}\\body dir $=2$};
\node[leaf] (X) at (6.7,-1.3) {\textbf{$X$ — send}\\no envelope};
\draw[edge] (root.south) -- node[lbl, midway, above left] {root} (ls.north);
\draw[edge] (root.south) -- node[lbl, midway, above right] {root} (rx.north);
\draw[edge] (ls.south) -- node[lbl, midway, left] {$h_L$} (L.north);
\draw[edge] (ls.south) -- node[lbl, midway, right] {$h_S$} (S.north);
\draw[edge] (rx.south) -- node[lbl, midway, left] {$h_R$} (R.north);
\draw[edge] (rx.south) -- node[lbl, midway, right] {$h_X$} (X.north);
\end{tikzpicture}
\caption{The four-leaf state output. The internal key is a nothing-up-my-sleeve
point; all leaves sit at depth two. $L$, $S$ and $R$ are the crossing leaves,
identical except for the direction byte in the body they push. $X$ is the
send leaf. There is no key-path spend.}
\label{fig:taproot}
\end{figure}

Every control block is 97 bytes:
\begin{equation}\label{eq:control}
\begin{aligned}
&[\mathit{parity}\mid\text{\code{0xc0}}]\;\|\;\text{\code{NUMS_XONLY}}\\
&\qquad\|\;\langle\text{sibling leaf hash}\rangle\;\|\;
\langle\text{sibling subtree hash}\rangle.
\end{aligned}
\end{equation}

\begin{table}[H]
\centering
\small
\caption{Control-block siblings by spent leaf.}
\label{tab:control}
\begin{tabular}{@{}lcc@{}}
\toprule
\textbf{Spent leaf} & \textbf{First sibling} & \textbf{Second sibling} \\
\midrule
$L$ (left)     & $h_S$ & $h_{RX}$ \\
$S$ (straight) & $h_L$ & $h_{RX}$ \\
$R$ (right)    & $h_X$ & $h_{LS}$ \\
$X$ (send)     & $h_R$ & $h_{LS}$ \\
\bottomrule
\end{tabular}
\end{table}

The three crossing leaves differ only in the direction byte of the body they
push (Section~\ref{sec:inscriptions}). Each leaf hash is therefore fixed when
the state output is created. Choosing a direction selects a leaf; it cannot
change the leaf's bytes. The current tree uses the
\code{FRG8CROSS}/\code{0x08} body. A state output whose sat equals
\code{GRANDFATHERED_SAT} may use the released
\code{FRG7CROSS}/\code{0x07} body and its own leaf hashes; the two versions
are never mixed in one tree.

\code{owner_xonly} is the 32-byte witness program of a BIP86 P2TR address
derived at
\begin{equation}
m/84'/\mathit{coin}'/0'/2/0,
\end{equation}
the ordinals branch. Its BIP86 even-$Y$ normalized output is the wallet's
current frog receive address and the state-leaf owner. The released
implementation tweaked the raw point without normalizing; a recovery scan
also watches that legacy address and derives the legacy state scripts. The
sender parses the recipient's P2TR address to obtain the key; a non-P2TR
destination is refused for custody.

\subsection{The cross leaf}

All three crossing leaves are the same script; only the last byte of the
pushed body (the direction) differs:
\begin{lstlisting}[language={},caption={Cross leaf template.},label={lst:cross-leaf}]
20 <owner:32>          ; PUSH32 owner
ad                     ; OP_CHECKSIGVERIFY
00                     ; OP_FALSE
63                     ; OP_IF
03 "ord"               ; PUSH 3 "ord"
01 01                  ; push tag 1
18 <content-type:24>   ; PUSH 24 "application/octet-stream"
01 03                  ; push tag 3
20 <parent:32>         ; PUSH32 mint txid (index 0)
00                     ; empty body tag
37 <body:55>           ; PUSH 55 bytes, last byte = dir; FRG8CROSS/0x08
68                     ; OP_ENDIF
51                     ; OP_TRUE
\end{lstlisting}

The first part is an ordinary \code{CHECKSIGVERIFY} against the owner key.
The second part is an ord envelope that is never executed as a condition:
\code{OP_FALSE} skips the \code{OP_IF} branch, so the envelope is data, not
control flow. The final \code{OP_TRUE} leaves a true stack value after the
signature check.

Every push is minimal. The content type is the 24-byte string
\code{application/octet-stream}, pushed with \code{0x18}. The parent tag is
the 32-byte mint txid for index $0$ --- the only reachable case, since a valid
mint carries exactly one envelope and it sits on input 0 --- or the txid plus
the minimal little-endian index (a push of length 33 to 36 bytes, opcode
\code{0x21} to \code{0x24}). The script is therefore \textbf{161 bytes} at index $0$, and up
to 165 bytes for a manually built nonzero index. The three direction leaves
are the same length; only the body's last byte differs.

The send leaf is deliberately smaller:
\begin{lstlisting}[language={},caption={Send leaf.},label={lst:send-leaf}]
20 <owner:32> ac
\end{lstlisting}
It is 34 bytes and contains no envelope. A transfer is therefore silent: it
moves the sat but publishes no state.

\subsection{Control blocks and witnesses}

A custody spend has three witness elements:

\begin{table}[H]
\centering
\small
\caption{Witness stacks for the two custody spends.}
\label{tab:witnesses}
\begin{tabularx}{\linewidth}{@{}p{0.16\linewidth}p{0.35\linewidth}X@{}}
\toprule
\textbf{Spend} & \textbf{\code{witness[0]}} & \textbf{\code{witness[1]}} / \textbf{\code{witness[2]}} \\
\midrule
hop & 64-byte BIP-340 \cite{bip340} signature $\|$ \code{0x01} (SIGHASH\_ALL) &
the chosen \code{dir} leaf script / 97-byte control block \\
transfer & 64-byte signature $\|$ \code{0x01} &
send leaf script / 97-byte control block \\
\bottomrule
\end{tabularx}
\end{table}

SPEC.md specifies that the signature is produced with the tweaked owner key
over the BIP-342 tapscript sighash of the leaf and the transaction's whole
prevouts and outputs. The reference implementation expresses that key as
\code{internal_key.tap_tweak(None)}. The SIGHASH\_ALL byte commits every
input and output, so a signature cannot be redirected to a different
transaction.

\subsection{The two-phase mint}

Ord needs an input to engrave an inscription, so a mint is a commit/reveal
pair. The commit output is a Taproot output with the BIP86 owner key
(\code{owner_xonly}) as internal key and one leaf:
\begin{lstlisting}[language={},caption={Commit output leaf.},label={lst:commit-leaf}]
<owner:32> OP_CHECKSIG OP_FALSE OP_IF <mint inscription envelope> OP_ENDIF
\end{lstlisting}
Its address is
$\mathrm{p2tr}(\mathit{owner},\mathrm{merkle\_root}(\mathrm{leaf}))$, and its
value is
\begin{equation}
1{,}000 + \mathrm{reveal\_fee\_budget}(\mathit{feerates}).
\end{equation}
The leaf demands the owner's signature. SIGHASH\_DEFAULT commits every
prevout and every output, so a copied witness cannot redirect the frog. A
legacy keyless commit, using \code{NUMS} as internal key with
\code{OP_TRUE || envelope}, remains revealable forever, but no new commit
may use it.

The reveal transaction then has the shape in Table~\ref{tab:reveal}.

\begin{table}[htbp]
\centering
\small
\caption{The mint reveal transaction.}
\label{tab:reveal}
\begin{tabularx}{\linewidth}{@{}p{0.12\linewidth}p{0.22\linewidth}X@{}}
\toprule
\textbf{Position} & \textbf{Role} & \textbf{Content} \\
\midrule
input 0 & commit spend & script path \code{[sig, leaf, control block]} \\
output 0 & frog & 1{,}000 sats to the wallet's frog P2TR key-path address \\
output 1 & mint message & OP\_RETURN \code{FRG 0x08 0x01 seed name} (\code{0x07} only for the grandfather) \\
output 2 & change & \code{m/84'/coin'/0'/1/index}, optional if $\ge$ dust \\
\bottomrule
\end{tabularx}
\end{table}

Every input uses sequence \code{0xFFFFFFFF}, so the transaction is not
replaceable. The anchor sat is the first sat of the commit output, which is
the first sat of the first funding input. Because the commit leaf embeds the
anchor sat, the commit address is derived after a probe coin-selection pass.

\subsection{Action transactions}

After the mint, a frog moves through the transactions in
Table~\ref{tab:actions}. All of them use version 2, lock time zero, and
sequence \code{0xFFFFFFFF} on every input. They are intentionally not
replaceable: a mint, prepare, hop or transfer may move an inscription or a
state card, and shaving an output to raise a fee can burn the anchor sat.

\begin{table}[htbp]
\centering
\small
\caption{The custody actions.}
\label{tab:actions}
\begin{tabularx}{\linewidth}{@{}p{0.12\linewidth}p{0.30\linewidth}X@{}}
\toprule
\textbf{Action} & \textbf{Input 0} & \textbf{Outputs} \\
\midrule
\textsc{prepare} & the frog output, P2TR key path with the tweaked frog key &
output 0 is the state output for $(\mathrm{sat},\mathit{mint},\mathit{owner})$;
outputs 1.. are funding change \\
\textsc{hop} & the state output, the leaf for the chosen direction &
output 0 is the next state output (same owner or a new P2TR owner);
outputs 1.. are funding change \\
\textsc{transfer} & the state output via the send leaf, or the mint output via
key path & output 0 is the recipient's ordinary address (any script, P2WPKH
included) \\
\bottomrule
\end{tabularx}
\end{table}

\begin{invariant}{One crossing per transaction}
A hop publishes exactly one inscription, in input 0. The anchor sat must
land in output 0 at offset 0. No funding input may carry a frog, and no
\code{FRG} envelope may appear anywhere in the crossing.
\end{invariant}

\textsc{prepare} exists only for a frog still in its mint or frog-address
output. The state leaf commits to the mint inscription id, which is created
by the mint reveal; an output of that transaction cannot contain its own
txid. The first state output is therefore created by either the first
transfer or the one-time \textsc{prepare}. After that, each crossing spends
the previous state output and creates the next one.

The player chooses $\mathit{dir}\in\{\text{left},\text{straight},\text{right}\}$
before signing; the transaction spends the matching leaf, so the body it
inscribes carries that \code{dir}. Choosing a direction is the whole of the
choice: the leaf's bytes are fixed by the state output, and the outcome is
decided by the block, not by the player. The wallet must show the chosen
direction in the preview and must not let it change after the preview is
approved.

A transfer deliberately pays the recipient's ordinary address, not a state
output. The state script commits to the mint id, which the recipient cannot
know until it has discovered the frog. The recipient re-enters custody with
\textsc{prepare} on its next action.

\subsection{Previews and the ordinary send}

A browser \textsc{send} (including \textbf{SEND ALL}), \textsc{hop},
\textsc{transfer} or \textsc{mint} is previewed before broadcast.
\textsc{send} and \textsc{transfer} display \code{PREVIEW TRANSACTION};
\textbf{SEND ALL} is titled \code{SEND ALL BITCOIN}; \textsc{mint} keeps
its \code{MINT!} button and opens the same review window; a hop is previewed
with its direction spelled out (\code{HOP LEFT},
\code{HOP STRAIGHT} or \code{HOP RIGHT}) and the leaf it will spend. The
preview shows the destination address and amount, change address and amount
(or ``one output'' for a transfer or \textbf{SEND ALL}), chosen fee rate and
actual fee in sats. A mint shows the commit and reveal txids, both change
outputs, both fees and the total size.

\code{CANCEL} performs no broadcast, persists no action and spends nothing.
Approval re-scans and refuses the preview if any input moved, an anchor's
ordinal ranges changed, or a coin gained a Hopper, then broadcasts the exact
previewed bytes. A mint approval broadcasts the commit and then the reveal
that spends it; if the reveal is rejected after the commit was accepted, the
error names the accepted commit and no mint action is recorded. A mint never
has an RBF path. Fee recommendations come from mempool.space
\code{/api/v1/fees/precise} (\code{fastestFee}, \code{halfHourFee},
\code{economyFee}); the chosen rate remains manually editable, and values
below the 1 sat/vB wallet relay floor are raised to it.

An ordinary wallet \textsc{send} is separate from these actions. Its inputs
use sequence \code{0xFFFFFFFD} (BIP125 opt-in), and the wallet may replace an
unconfirmed send only under a narrow policy: the original is a wallet-created
plain BIP84 send; the backend returns exactly the recorded serialization;
every original input is a plain, range-known, non-Hopper BIP84 coin and the
current view still contains no Hopper at those outpoints or sats; every
original output keeps its exact script and value; the replacement appends
confirmed plain BIP84 funding and one new change output; and it pays a
strictly higher absolute fee and feerate, including BIP125's incremental
relay margin. The wallet never replaces a HOPPER action. The transaction
builder recomputes the FIFO sat ranges and refuses a replacement if a known
Hopper sat would land in its fee.

\textbf{SEND ALL.} A SEND variant sweeps the plain-BTC balance into exactly
one recipient output. It selects every eligible plain BIP84 coin: complete,
value-matched ranges, no Hopper or state output, and no protected sat;
unconfirmed coins are included only when the wallet permits them. There is no
change output, and the actual fee is deducted from the total, so the recipient
receives $\mathit{totalIn}-\mathit{actualFee}$. The build refuses when the
fee consumes the balance or leaves the recipient below the P2WPKH dust limit,
and approval repeats the same exact-bytes and input checks as SEND. The
preview is titled \code{SEND ALL BITCOIN} and states that there is one output
and no change.

\subsection{Sat assignment}

Ordinals move FIFO over inputs in order. The protocol depends on three
consequences.

\begin{enumerate}[leftmargin=1.6em,itemsep=0.22em]
  \item The anchor sat is the first sat of input 0. For a mint, that is the
        first sat of the commit output.
  \item After a spend, the anchor sat is at offset 0 of output 0 if and only
        if output 0 receives value and no earlier output claims sats first.
  \item A fee is taken from the tail of the frog's range or from funding; it
        is never taken from the anchor's position.
\end{enumerate}

The wallet and the indexer both verify these properties with ordinal range
arithmetic. A coin whose ranges cannot be attributed is not spent.

\subsection{Addresses, discovery and history}

\begin{itemize}[leftmargin=1.4em,itemsep=0.3em]
  \item The mint pays the wallet's \textbf{frog address}: the BIP86 P2TR
        output at \code{.../2/0}. It never pays a BIP-84 receive address.
        Coins at the frog address are never funding candidates.
  \item The BIP-84 receive script is \textbf{fixed at \code{.../0/0}} in
        v8. This is a deliberate simplicity and recovery tradeoff, not an
        accidental gap: the watch set is fixed, paginated to exhaustion, and
        a restored wallet reuses the same script. Address rotation and
        privacy are future wallet features, not a consensus rule. The UI
        must not claim that a receive address retires.
  \item The state output's script is per
        $(\mathrm{sat},\mathit{mint\ id},\mathit{owner})$. It matches no
        derivation path. A wallet finds it by deriving the script from the
        frog record and by watching its frog address.
  \item A frog sent by a peer arrives at the ordinary address and is
        discovered by the wallet's address scan plus the indexer or ord. It
        is imported by reading its mint inscription and replaying its history
        from the mint outpoint.
  \item History is wallet-scoped, not a raw indexer feed. It merges local
        actions, events whose Hopper sat belongs to this wallet, settled
        outputs observed on its watched scripts, internal change, and mempool
        rows. A confirmed ordinary deposit or a Hopper output is retained as
        a \code{receive} row, so it stays visible after the coin is spent.
        Change from our own transactions is retained separately as a
        \code{CHANGE} row; it is never presented as an incoming
        \code{RECEIVE}. Events for unrelated public or NPC Hoppers are
        filtered out.
  \item A \code{MINT} row does not depend only on the browser that broadcast
        it. When the owned mint output is still at its mint outpoint, the
        wallet persists a chain-derived mint row, so a restored seed shows
        the same \code{MINT} and \code{HOPPER ARRIVED} rows without a local
        action ledger. A Hopper received from a peer is not at its mint
        outpoint and does not create a \code{MINT} row.
  \item Rows for one txid share its best-known height and sort as a group.
        The feed is newest first, so inside a group the newest row is on top:
        \code{HOPPER ARRIVED} $\to$ \code{CHANGE} $\to$ \code{MINT}, and
        \code{CHANGE} $\to$ \code{SEND}. A mint's commit change is grouped
        with the reveal's \code{HOPPER ARRIVED} row by the minted sat and by
        both txids. Older actions with no stored sat recover it from the
        arrival's frog/reveal record; a missing commit is read once from
        input 0 of the reveal and persisted. A missing sat is never treated
        as sat 0. Grouping is ultimately by reveal/commit txid cluster, so it
        works even when no row stores a sat. Shortened txids and addresses
        are copy buttons: the visible text is abbreviated, but the copied
        text is the complete value.
  \item The first completed refresh after the history-ledger upgrade performs
        two one-time, versioned address-history backfills: receive addresses
        the wallet handed out, and change addresses it has used. The receive
        pass recovers deposits whose coins were spent before the receive
        ledger existed and skips outgoing action/commit txids; the change
        pass recovers spent internal change and migrates older ledgers that
        had stored it as a \code{RECEIVE}. \code{RESCAN} clears both markers
        so recovery can be forced.
\end{itemize}

\begin{invariant}{No key-path escape hatch}
A state output is a UTXO spent exactly once. The body binds the sat, the mint
and the direction, so a witness cannot be replayed against another frog. The
NUMS internal key ensures that the only ways out are the four leaves.
\end{invariant}
